% SPDX-License-Identifier: CC-BY-4.0
% Copyright (c) 2026 Martin Schröder <info@swedishembedded.com>

\section{Lessons: a formal account}\label{sec:formalism}

Definitions follow the implementation (\code{core::speech\_lesson}); each statement below is the
specification of a test in that module.

\paragraph{Directives.}
A transcript $u$ is read to at most one directive. Let $\operatorname{norm}(u)$ be $u$ with leading and
trailing white space removed and letters lower-cased. $u$ is an \emph{imitation} if $\operatorname{norm}(u)$
begins with one of the openers ``talk as follows'', ``speak as follows'', ``say it like this'', ``say this
like this'', ``say it as follows''; its example is the remainder of $u$ after the opener, with leading
separators (colon, comma, full stop, semicolon, hyphen, space) removed, and it must be non-empty. $u$ is a
\emph{respelling} if it begins with ``pronounce '' or ``say '' and the remainder contains a first ``\,as\,'' that
splits it into a word $w$ and a spelling $s$ (quotes and a trailing full stop or exclamation mark removed),
both non-empty and neither containing a space. Any other $u$ is not a directive. The reading is
deliberately conservative: a false lesson is worse than a missed one, and a missed one costs the user
a repetition.

\paragraph{Lessons and objectives.}
A lesson is a record $(h, e, o, \pi)$: what was heard $h$ (kept as heard), an exemplar $e$, an objective
$o\in\{\textsc{Lexicon},\textsc{InContext},\textsc{Supervised},\textsc{Preference}\}$ and the portrayal $\pi$ of the voice it
applies to, which is never blank. A respelling $(w,s)$ yields the single lesson $(h, (w{\to}s), \textsc{Lexicon}, \pi)$.
An imitation with example text $x$ yields, only when the user's example was recorded (digest $d$ of its audio), the
three lessons $(h,(d,x),o,\pi)$ for $o\in\{\textsc{InContext},\textsc{Supervised},\textsc{Preference}\}$;
without the recording it yields none and says why, because the words alone do not carry how they were said.
A lesson names audio by digest and never holds it.

\paragraph{Lexicon.}
A lexicon $L$ maps lower-cased words to spellings; \emph{teach}$(w,s)$ sets $L(\operatorname{lower}(w))=s$, so a later lesson for
a word replaces an earlier one. Applying $L$ to a text $t$ splits $t$ into maximal runs of letters, digits
and apostrophes (words) and everything else; each word $v$ with $\operatorname{lower}(v)\in L$ is replaced by
$L(\operatorname{lower}(v))$, all else is kept, including punctuation and spacing. Consequently:
(i) the empty lexicon is the identity; (ii) application is a single pass, never recursive; (iii) a possessive such as
``Jefferson's'' is a different word from ``Jefferson'' and is untouched; (iv) application changes only the text handed to
the voice, never the text the persona wrote, and never the speaker profile or its portrayal.

\paragraph{Reward and comparison of takes.}
A take has a judged naturalness $n\in[0,1]$ and a word error rate $e\in[0,1]$ (the share of its words lost when it is heard
again). Its reward is $r=n-\lambda e$ with $\lambda=5$ by default: losing a fifth of the words costs more than the most natural
delivery can earn back. For a group of $k\ge2$ takes of one sentence with rewards $r_i$, mean $\bar r$ and standard deviation
$\sigma$, the advantage is $A_i=(r_i-\bar r)/\sigma$, and $A_i=0$ for all $i$ when $\sigma<10^{-6}$ or $k<2$: a tie teaches
nothing. The preference pair of a group is $(\arg\max r,\arg\min r)$ when these differ. Advantages are centred and invariant
to adding a constant to the rewards, so a long sentence is compared with other takes of itself, not with a short one.

\paragraph{Planned objective.}
For the \textsc{Preference} objective the policy is the synthesizer's first-codebook token distribution,
optimised with the clipped group-relative surrogate of \citet{shao2024deepseekmath} with a KL term to the frozen
original; the remaining codebooks and the codec stay frozen. This is specified and not yet implemented
(\cref{sec:discussion}).
